\begin{center}\begin{tikzpicture}[font=\small,>=stealth,line join=round]
\draw[gray!40,rounded corners=3pt] (0,0) rectangle (14.2,8.5);
\node at (7.1,7.9) {例2：不必一直递增，奇函数只要求对称伙伴};
\begin{scope}[shift={(3.6,4.2)},x=1.15cm,y=0.85cm]\draw[->,gray] (-2.2,0)--(2.4000000000000004,0) node[right]{$x$};\draw[->,gray] (0,-3)--(0,3.2) node[left]{$y$};\begin{scope}\clip (-2.2,-3) rectangle (2.2,3);\draw[blue,thick,domain=-2.2:2.2,samples=120] plot(\x,{\x*\x*\x-2*\x});\end{scope}\draw[orange!85!black,dashed] (-1,1)--(1,-1);\fill[orange!85!black] (-1,1) circle(2pt) node[above left]{$(-1,1)$};\fill[orange!85!black] (1,-1) circle(2pt) node[below right]{$(1,-1)$};\end{scope}\node[align=left,anchor=west] at (7.8,4.7) {$f(x)=x^3-2x$\\[12pt]定义域 $\mathbb R$ 成对\\[12pt]$f(-x)=-f(x)$\\[12pt]任意点绕原点半转后仍在曲线上};
\node[text width=12.8cm,align=center] at (7.1,.65) {橙色点对只是示例；正文的恒等式才证明对每个输入都成立。};
\end{tikzpicture}\end{center}
